\section{Mistral Large 4（Le Chonk）公開プレビュー}
Mistral AIは、総パラメータ約1T・アクティブ49BのネイティブマルチモーダルモデルMistral Large 4を告知した。APIは即日提供、open weightsは10月末予定としている。サイバー防御やvisual groundingを強調するが、公開された評価値は同社側の主張であり、独立検証ではない。
\dailyxsource{Mistral AI (@MistralAI)}{https://x.com/MistralAI/status/2107457414387622310}

\section{Cline、Large 4の比較結果を投稿}
Clineは、Mistral Large 4がセキュリティ関連評価で他のフロンティアモデルを上回った背景として、課題への応答方針の差を挙げた。ベンチ名、課題セット、評価定義、再現手順は今回未確認であり、Cline側の比較主張として扱う。
\dailyxsource{Cline (@cline)}{https://x.com/cline/status/2107561157787824347}

\section{OpenAI Decisions API、公開ベータ開始}
OpenAI DevelopersはDecisions APIを全開発者向け公開ベータとして告知した。モデル・ツール・アクションの選択を近リアルタイムで行う用途で、GPT-6 Lunaを使いpredicate / choice / scoreを返す。Grok intakeでは入力課金のみで、Responses API経由より最大10倍速いという公式説明も確認されている。
\dailyxsource{OpenAI Developers (@OpenAIDevs)}{https://x.com/OpenAIDevs/status/2107573382229188645}

\section{OpenAI、Codex auto-reviewとAPI usage tierを更新}
OpenAIのTiboはDevDay Day 2相当の更新として、Approve for me、builders向けAPI簡素化、meeting notes、Decisions APIを列挙した。Grok intakeではusage tierが5段階からBuild / Launch / Growの3段階へ整理されたことを公式changelogで確認。一方、auto-reviewの対象条件や「2--10\%」の説明は一次仕様未確認である。
\dailyxsource{Tibo (@thsottiaux)}{https://x.com/thsottiaux/status/2107575657014468879}

\section{Google、Nano Banana 2.1を告知}
Googleは画像生成・編集モデルNano Banana 2.1を告知した。前モデル比でvisual design、mask-based editing、subject consistencyが向上したとしている。Arena.aiも初期順位を投稿したが、パラメータ数、APIモデルID、料金などはGrok収集時点で公式資料を確認できていない。
\dailyxsource{Google (@Google)}{https://x.com/Google/status/2107501209154204148}

\section{Google DeepMind、EmbeddingGemma 2を公開}
Google DeepMindは、テキスト・code・image・audio・videoを共有空間へ写すオンデバイス向け埋め込みモデルEmbeddingGemma 2を公開した。Grok intakeでは740Mパラメータ、Apache 2.0、Gemma 4系で、Hugging FaceとKaggleにweightsを提供すると確認されている。
\dailyxsource{Google DeepMind (@GoogleDeepMind)}{https://x.com/GoogleDeepMind/status/2107502286758895878}

\section{Anthropic、Cyber Verification Programを拡大}
AnthropicはCyber Verification Programを拡大し、検証済みセキュリティ専門家にMythos 5.1、Opus 5.5、Sonnet 5.5へのアクセスを広げると発表した。新しい専門家向けティアも設けるとしているが、申請資格や監査要件は投稿本文では未確認である。
\dailyxsource{Anthropic (@AnthropicAI)}{https://x.com/AnthropicAI/status/2107546569654636883}

\section{Google Earth AI、公衆衛生での地理空間推論を紹介}
Googleは、衛星・人口動態・地理空間モデルを組み合わせたGoogle Earth AI研究を投稿した。コンゴ民主共和国のエボラ対応でWHO AFROがGeospatial Reasoning agentのプロトタイプを使い、数分で48の曝露集落と45,500人超を特定したとする。数値は公式スレッドのパートナー事例で、独立検証は行われていない。
\dailyxsource{Google (@Google)}{https://x.com/Google/status/2107578315527623122}

\section{Reflection AI Beam、窓内では続報が継続}
窓内の二次投稿は、Reflection AIのopen-weightモデルBeamについて、501B総量・23B activeのMoE、コーディング・推論・エージェント用途、weightsは今月後半予定などを伝えた。初報として参照されるReuters記事は観測窓開始前で、Reflection公式X投稿も今回未確認のため、続報・議論として扱う。
\dailyxsource{AI Tech Updates (@AIInfoShare)}{https://x.com/AIInfoShare/status/2107484640592400499}

\section{OpenAI textGrain、EU向け出自表示の話題が継続}
窓内の投稿は、OpenAIがEU AI Act対応としてChatGPTとCodexのテキスト出力へ統計的ウォーターマークtextGrainを導入すると伝えた。グローバルAPIは対象モデルでオプトイン、検出器は承認された研究者・専門家向けとされるが、今回のGrok検索ではOpenAI公式X投稿を確認できていない。
\dailyxsource{まっすー (@Massu15)}{https://x.com/Massu15/status/2107588418888958431}

\section{micro1、PersonalAgentBenchを公開}
micro1のCEOは、個人エージェントの信頼性を測るPersonalAgentBenchを公開したと投稿した。Instinct、Muse、Grok Bot、Gemini Sparkを比較し、正しい情報へ到達しても過剰共有や創作が残るとする。課題定義、勝敗表、報酬条件は投稿本文だけでは未確認である。
\dailyxsource{Ali Ansari (@aliansarinik)}{https://x.com/aliansarinik/status/2107497652871155953}

\section{GitHub、コードレビュー用ReviewBenchを告知}
GitHubはAIコードレビューエージェント向けオープンベンチReviewBenchを告知した。公式投稿では1億390万件のpull request分析に基づき、19言語・219 PRで構成するとしている。データセット、評価ハーネス、提出先URLは投稿本文からは確認できていない。
\dailyxsource{GitHub (@github)}{https://x.com/github/status/2107540144949473754}

\section{GitHub、agent-scaleに向けGit基盤を再構築}
GitHubは、2025年9月から2026年8月の間にGit活動が月間473.3B eventsまで倍以上になったとし、サービスを止めずにGit基盤を再構築してagent-scaleのソフトウェア開発に備えると投稿した。イベントの定義や詳細はリンク先X Articleに委ねられ、Grok intakeでは本文未取得である。
\dailyxsource{GitHub (@github)}{https://x.com/github/status/2107589127051063771}

\section{OpenCode、Mistral Large 4を即日提供}
OpenCodeはMistral Large 4を同ツールで利用可能にし、今後2週間50\%オフと投稿した。Mistral側の「API today」という告知と時刻は整合するが、対象プラン、通常価格、割引の適用範囲は投稿本文にない。Large 4のエージェント製品への即日搭載例として記録する。
\dailyxsource{OpenCode (@opencode)}{https://x.com/opencode/status/2107591389592867155}
